\clearpage
\VerdierPhysicalPageAnchor{0129}%
\begingroup
\small
\setlength{\emergencystretch}{1.5em}%
\setlength{\abovedisplayskip}{3pt plus 2pt minus 2pt}%
\setlength{\belowdisplayskip}{3pt plus 2pt minus 2pt}%

\VerdierSemanticAnchor{ch2-2-1-1-sm3-conditions}
\VerdierSemanticAnchor{ch2-2-1-1-sm3-i}
\noindent\textit{i)} Il existe un morphisme $s$ de $S$, de source $Y$, tel que
$sf=sg$.

\VerdierSemanticAnchor{ch2-2-1-1-sm3-ii}
\noindent\textit{ii)} Il existe un morphisme $t$ de $S$, de but $X$, tel que
$ft=gt$.

\VerdierSemanticAnchor{ch2-2-1-1-saturated}
\noindent Un système multiplicatif est dit \emph{saturé} s'il possède la
propriété suivante :

\VerdierSemanticAnchor{ch2-2-1-1-sm4}
\VerdierSemanticAnchor{ch2-2-1-1-saturated-sm4}
\smallskip
\noindent\textbf{SM4)} Si $f,g$ et $h$ sont trois morphismes composables et si
les morphismes $fg$ et $gh$ appartiennent à $S$, alors le morphisme $g$
appartient à $S$.

\VerdierSemanticAnchor{ch2-2-1-2}
\medskip
\noindent\textbf{2.1.2.}\quad
\VerdierSemanticAnchor{ch2-2-1-2-compatible-triangulation}
\VerdierSemanticAnchor{ch2-2-1-2-compatible}
Soit $\mathcal D$ une catégorie triangulée. Un système multiplicatif $S$ de
$\mathcal D$ est dit \emph{compatible avec la triangulation} s'il possède les
propriétés suivantes :

\VerdierSemanticAnchor{ch2-2-1-2-sm5}
\smallskip
\noindent\textbf{SM5)} Pour tout élément $s$ de $S$ et tout entier
$n\in\mathbb Z$, le morphisme $s[n]$ appartient à $S$. Autrement dit, $S$ est
\emph{stable par les automorphismes de translation}.

\VerdierSemanticAnchor{ch2-2-1-2-sm6}
\smallskip
\noindent\textbf{SM6)} Pour tout couple de triangles distingués :

\VerdierSemanticAnchor{ch2-2-1-2-sm6-triangle-pair}
\[
  (X,Y,Z,u,v,w), \qquad (X',Y',Z',u',v',w')
\]
\noindent et tout diagramme commutatif :

\VerdierSemanticAnchor{ch2-2-1-2-sm6-square}
\begin{center}
  \begin{tikzcd}[column sep=4.2em,row sep=2.5em]
    X \arrow[r,"u"] \arrow[d,"s"'] & Y \arrow[d,"s'"] \\
    X' \arrow[r,"u'"'] & Y'
  \end{tikzcd}
  \quad ,
\end{center}
\VerdierSemanticAnchor{ch2-2-1-2-sm6-triangle-morphism}
\VerdierSemanticAnchor{ch2-2-1-2-sm6-completion}
\noindent où $s$ et $s'$ sont des éléments de $S$, il existe un morphisme
appartenant à $S$ : $s''\colon Z\to Z'$ tel que $(s,s',s'')$ soit un morphisme
de triangles.

\VerdierSemanticAnchor{ch2-2-1-3}
\medskip
\noindent\textbf{Remarque 2.1.3.}\quad
\VerdierSemanticAnchor{ch2-2-1-3-item1}
\VerdierSemanticAnchor{ch2-2-1-3-1}
\textbf{1)} Soient $\mathcal D$ une catégorie triangulée, $S$ un système
multiplicatif de $\mathcal D$ compatible avec la triangulation et
$\mathcal D^\circ$ la catégorie triangulée opposée à $\mathcal D$. L'ensemble
$S$ de morphismes de $\mathcal D^\circ$ est un système multiplicatif de
$\mathcal D^\circ$ compatible avec la triangulation de $\mathcal D^\circ$
(1.1.7).

\VerdierSemanticAnchor{ch2-2-1-3-item2}
\VerdierSemanticAnchor{ch2-2-1-3-2}
\smallskip
\noindent\textbf{2)} Soient $\mathcal D$ une catégorie triangulée et $S$ un
ensemble de morphismes de $\mathcal D$ possédant les propriétés SM1), SM5) et
SM6). Alors l'ensemble $S$ possède aussi la propriété SM2).
\VerdierSemanticAnchor{ch2-2-1-3-2-sm2}
En effet, il suffit de montrer que tout diagramme :

\VerdierSemanticAnchor{ch2-2-1-3-item2-input-boundary}
\VerdierSemanticAnchor{ch2-2-1-3-2-input-boundary}
\begin{center}
  \begin{tikzcd}[column sep=4.0em,row sep=2.6em]
    & X' \arrow[d,"s\in S"] \\
    Z \arrow[r,"f"] & X
  \end{tikzcd}
\end{center}
\endgroup
